\subsection{Caracteres, bits y números}
\label{alg:2-8}

\subsubsection{ord() y chr()}

Convierten entre un carácter y su código. Es la forma de hacer aritmética con
letras: pasar de letra a índice, desplazar un cifrado, contar por alfabeto.

\begin{lstlisting}[language=Python]
ord("a"), chr(98)                  # (97, 'b')
c = "d"
i = ord(c) - ord("a")              # letra -> indice 0..25: 3
chr(ord("a") + (i + 2) % 26)       # desplazar 2 letras: 'f'
cnt = [0] * 26
for ch in "banana":
    cnt[ord(ch) - 97] += 1         # cnt[0] == 3 (tres 'a')
\end{lstlisting}

Al leer con \texttt{sys.stdin.buffer} no hace falta \texttt{ord}: indexar un
\texttt{bytes} ya da el código (\texttt{b"a"[0] == 97}).

\begin{lstlisting}[language=Python]
import string
string.ascii_lowercase             # 'abcdefghijklmnopqrstuvwxyz'
string.ascii_uppercase             # 'ABC...Z'
string.digits                      # '0123456789'
\end{lstlisting}

\vspace{0.3cm}
\subsubsection{Operaciones de bits}

\begin{lstlisting}[language=Python]
x = 12                             # 0b1100
bin(x), bin(x)[2:]                 # ('0b1100', '1100')
format(5, "08b")                   # '00000101': 8 bits con ceros
int("1010", 2)                     # 10
x.bit_length()                     # 4: bits necesarios
bin(x).count("1")                  # 2 bits en 1 (x.bit_count() en 3.10+)
x & -x                             # 4: bit encendido mas bajo
x & (x - 1)                        # 8: apaga ese bit
j = 1
(x >> j) & 1                       # leer el bit j: 0
x | (1 << j)                       # encender el bit j: 14
x & ~(1 << j)                      # apagar el bit j
x ^ (1 << j)                       # invertir el bit j: 14
\end{lstlisting}

Los enteros de Python no tienen ancho fijo, así que \texttt{\textasciitilde x} vale
$-x-1$ en vez de invertir 32 o 64 bits; para complementar solo $n$ bits se usa
\texttt{x \^{} ((1 <{}< n) - 1)}.

Recorrer todas las submáscaras de una máscara, de mayor a menor:

\begin{lstlisting}[language=Python]
mask = 0b1011
sub = mask
while sub:
    ...                            # procesar sub: 1011, 1010, 1001, 1000, 11, 10, 1
    sub = (sub - 1) & mask
# la submascara vacia (0) no entra en el ciclo: procesarla aparte si hace falta
\end{lstlisting}

Hacer esto para todas las máscaras de $n$ bits cuesta \bigO{\pow{3}{n}}, no
\bigO{\pow{4}{n}}: es el límite usual de los DP sobre subconjuntos de subconjuntos.

\vspace{0.3cm}
\subsubsection{Enteros grandes}

Los enteros de Python no se desbordan: factoriales, productos de números de
\ten{18} o potencias enormes se calculan exactos, sin \texttt{long long} ni módulo
forzado. Dos advertencias. Operar con números de miles de dígitos ya no cuesta
\bigO{1}. Y desde Python 3.11 (y desde los parches de seguridad 3.7.14, 3.8.14, 3.9.14 y 3.10.7), convertir entre entero y texto un número de más de
4300 dígitos lanza \texttt{ValueError}, tanto al imprimirlo como al leerlo con
\texttt{int()}:

\begin{lstlisting}[language=Python]
import sys, math
if hasattr(sys, "set_int_max_str_digits"):   # no existe si no hay limite
    sys.set_int_max_str_digits(0)            # 0 = sin limite
print(math.factorial(2000))                  # 5736 digitos
\end{lstlisting}

El \texttt{hasattr} hace que la misma línea sirva en cualquier versión: en las
versiones sin el límite, la función no existe.

\vspace{0.3cm}
\subsubsection{Números reales}

\begin{lstlisting}[language=Python]
import math
INF = float("inf")                 # mayor que cualquier numero (o math.inf)
0.1 + 0.2 == 0.3                   # False
EPS = 1e-9
abs((0.1 + 0.2) - 0.3) < EPS       # True: comparar con tolerancia
math.isclose(1e-12, 0)             # False!
math.isclose(1e-12, 0, abs_tol=EPS)  # True
\end{lstlisting}

\texttt{math.isclose} usa solo tolerancia relativa por defecto, y cerca de cero eso
no sirve: hay que pasarle \texttt{abs\_tol}. Siempre que se pueda, mejor evitar los
reales: \texttt{a/b < c/d} se compara como \texttt{a*d < c*b} (con $b, d > 0$) y
queda exacto.
